\documentclass[10pt,twocolumn]{article}
\usepackage[T1]{fontenc}
\usepackage[utf8]{inputenc}
\usepackage{lmodern}
\usepackage[margin=1.8cm,top=2.4cm,bottom=2cm,columnsep=0.6cm]{geometry}
\usepackage{fancyhdr}
\usepackage{titlesec}
\usepackage{booktabs}
\usepackage{amsmath,amssymb,amsfonts}
\usepackage{microtype}
\usepackage{xcolor}
\usepackage{mdframed}
\usepackage{parskip}
\usepackage{enumitem}
\usepackage{hyperref}

\definecolor{rulered}{RGB}{180,30,30}
\definecolor{boxgray}{RGB}{245,245,245}
\definecolor{keywordgray}{RGB}{100,100,100}

\pagestyle{fancy}
\fancyhf{}
\fancyhead[L]{\textsc{\small Claude Mag}}
\fancyhead[C]{\small\textit{Machine Learning Research Digest}}
\fancyhead[R]{\small 2026-10-07}
\fancyfoot[C]{\small\thepage}
\renewcommand{\headrulewidth}{0.8pt}

\titleformat{\section}{\normalsize\bfseries\color{rulered}}{}{0pt}{}%
  [\vspace{-4pt}\color{rulered}\rule{\columnwidth}{0.4pt}]
\titlespacing{\section}{0pt}{8pt}{4pt}

\hypersetup{colorlinks=true,urlcolor=rulered,linkcolor=black}

\begin{document}
\setlength{\parindent}{0pt}
\setlength{\parskip}{4pt}

{\small\color{keywordgray}\textbf{Keywords:}
  data selection \textbullet{}
  post-training \textbullet{}
  gradient alignment \textbullet{}
  LLM fine-tuning}\par\vspace{4pt}

{\LARGE\bfseries\raggedright
  LESSER: Post-Training Data Selection with Output-Layer Gradients}\par\vspace{4pt}

{\large\itshape\raggedright
  Output-Layer Gradients Are a 9.7$\times$ Cheaper Proxy for Full-Gradient Data
  Selection While Preserving Batch-Level Alignment and Downstream Performance}\par\vspace{6pt}

{\small\textbf{By} Lyuxin David Zhang, Eric Wong, Surbhi Goel \& Anton Xue
  (University of Pennsylvania)}\par\vspace{2pt}

{\small\color{keywordgray}arXiv:2610.03702 \textbullet{} October 2026}
\par\vspace{2pt}\noindent{\color{rulered}\rule{\columnwidth}{0.6pt}}\vspace{6pt}

Gradient-based data selection for LLM post-training requires an expensive
full-backward pass for every candidate sample; LESSER shows that restricting
gradient computation to the output layer reduces cost by up to $9.7\times$
while matching full-gradient performance, because batches selected by both
methods remain well-aligned with the validation gradient even when
individual sample rankings differ.

\begin{mdframed}[backgroundcolor=boxgray,linecolor=rulered,linewidth=1pt,
    innerleftmargin=6pt,innerrightmargin=6pt,
    innertopmargin=6pt,innerbottommargin=6pt]
\textbf{\small AT A GLANCE}\par\vspace{4pt}
\begin{description}[leftmargin=0pt,labelsep=4pt,itemsep=2pt,parsep=0pt,
    font=\small\bfseries]
  \item[Problem:] Full-parameter gradient computation for post-training data
    selection is intractable at scale.
  \item[Why it matters:] The right training data can substantially improve
    downstream performance; efficient selection unlocks practical use.
  \item[Method:] Restrict gradient computation to the output (language-model
    head) layer only; use as a drop-in wrapper for any selection algorithm.
  \item[Result:] 9.7$\times$ FLOP reduction for SFT and 3.0$\times$ for RL
    selection, with no downstream performance loss.
  \item[Evidence:] Evaluated across multiple gradient-based selection methods
    on SFT and RL post-training benchmarks.
\end{description}
\end{mdframed}

\section{The Big Picture}

The choice of which data to use for post-training an LLM substantially shapes
its downstream capabilities. Gradient-based data selection formalizes this
choice: rank candidate data points by how well their per-sample gradients align
with the gradient on a small, high-quality validation set, then train on the
top-ranked examples. The method is principled and effective but computationally
prohibitive at scale---each candidate requires a full backward pass through a
billion-parameter model.

LESSER (Last-layer Efficient Sample Selection via Output-layer gRadients) asks:
does the output layer's gradient contain enough information to perform
effective selection, at a fraction of the cost of backpropagating through
the entire network?

\section{Why Existing Approaches Fall Short}

\textbf{Full-gradient cost is prohibitive.} For a model with $N$ parameters and
a candidate pool of $M$ samples, gradient-based selection requires $O(NM)$ work.
At $N = 7 \times 10^9$ and $M = 10^5$, this is impractical without
approximation.

\textbf{Embedding-layer proxies introduce extra hyperparameters.} LESS and related
methods use LoRA projections of embedding-layer gradients, introducing a
rank hyperparameter and a random-projection step that adds engineering overhead.

\textbf{Sample-level alignment is not the operative quantity.} Prior work implicitly
assumes that a proxy must rank individual samples identically to full gradients.
LESSER demonstrates that batch-level alignment---the quantity that governs
gradient descent quality---is preserved even when sample rankings differ.

\section{Core Mechanism}

LESSER restricts gradient computation to the \textbf{output layer} (language
model head)---the final linear projection $W_\text{head}$ from hidden states
to vocabulary logits. This layer has far fewer parameters than the full model
but its gradient directly reflects the training objective.

The output-layer gradient for a data point $x$ is:
\[
\nabla_{W_\text{head}} \ell(x, \theta) = \frac{\partial \ell}{\partial \text{logits}} \cdot h_x^\top
\]
where $h_x$ is the final hidden state and $\frac{\partial \ell}{\partial \text{logits}}$
is the softmax loss gradient. This requires only a single matrix multiplication
and does not backpropagate through transformer layers.

\section{Technical Formulation}

Standard gradient-based selection scores candidate $x$ as:
\[
s(x) = \left\langle \nabla_\theta \ell(x, \theta),\; \nabla_\theta \mathcal{L}(\mathcal{V}, \theta) \right\rangle
\]

LESSER replaces $\theta$ with the output-layer parameters $\theta_L$:
\[
s_\text{LESSER}(x) = \left\langle \nabla_{\theta_L} \ell(x, \theta),\; \nabla_{\theta_L} \mathcal{L}(\mathcal{V}, \theta) \right\rangle
\]

The key empirical finding---that batch-level alignment is preserved---can be
stated as follows. Let $B_\text{full}$ and $B_\text{LESSER}$ be the top-$k$
batches selected by each method. Then empirically:
\[
\left\langle \nabla_\theta \mathcal{L}(B_\text{LESSER}, \theta),\;
\nabla_\theta \mathcal{L}(\mathcal{V}, \theta) \right\rangle
\approx
\left\langle \nabla_\theta \mathcal{L}(B_\text{full}, \theta),\;
\nabla_\theta \mathcal{L}(\mathcal{V}, \theta) \right\rangle
\]

The intuition is that the last layer aggregates information from all earlier
layers, so its gradient captures the dominant component of the full gradient
along the data manifold.

\section{Learning or Inference Procedure}

\begin{enumerate}[itemsep=2pt,parsep=0pt]
  \item \textbf{Cache validation gradient:} Run a forward-backward pass on the
    validation set using only output-layer parameters. Cache
    $g_\mathcal{V} = \nabla_{\theta_L} \mathcal{L}(\mathcal{V}, \theta)$.
  \item \textbf{Score candidates:} For each candidate $x_i$, run a forward
    pass; compute $g_i = \nabla_{\theta_L} \ell(x_i, \theta)$ via one matrix
    multiplication; score $s_i = \langle g_i, g_\mathcal{V} \rangle$.
  \item \textbf{Select top-$k$:} Rank candidates by $s_i$; select the top $k$.
  \item \textbf{Post-train:} Fine-tune on selected data with standard SFT or
    RL objective.
\end{enumerate}

\section{What the Guarantee Says}

There is no formal guarantee that output-layer ranking preserves full-gradient
ranking at the sample level---and indeed it does not, in general. The operative
guarantee is empirical and at the batch level: across all tested settings,
batches selected by LESSER achieve gradient alignment with the validation set
indistinguishable from those selected by full-gradient methods, and deliver
matching downstream task performance.

\section{Experimental Findings}

\begin{itemize}[itemsep=2pt,parsep=0pt]
  \item \textbf{FLOP reduction:} LESSER reduces feature-extraction cost by
    \textbf{9.7$\times$} for SFT and \textbf{3.0$\times$} for RL benchmarks.
  \item \textbf{Performance parity:} On downstream evaluations, LESSER matches
    full-gradient selection within measurement noise across all settings.
  \item \textbf{Sample vs.\ batch alignment:} Output-layer and full-gradient
    Spearman sample-ranking correlations vary widely, but batch-level gradient
    alignment with the validation set is consistently preserved.
  \item \textbf{Compatibility:} LESSER wraps LESS, DataInf, and other
    selection algorithms without modifying the selection logic.
\end{itemize}

\section{Ablations and Interpretation}

\textbf{Layer depth sweep:} Ablations over which layer's gradients to use show
a monotone improvement with depth; the output layer is best. Mid-layer gradients
reduce cost but widen the quality gap.

\textbf{Validation set size:} Performance degrades gracefully as validation size
shrinks; 128 examples suffice for reliable selection.

\textbf{SFT vs.\ RL gap:} The FLOP saving is larger for SFT (9.7$\times$) than
RL (3.0$\times$) because RL requires gradients over longer generated sequences,
making each full backward pass more expensive and the relative saving smaller.

\section*{Reference}

Lyuxin David Zhang, Eric Wong, Surbhi Goel \& Anton Xue.
\textbf{LESSER: Post-Training Data Selection with Output-Layer Gradients}.
arXiv:2610.03702, October 2026.\\
\url{https://arxiv.org/abs/2610.03702}

\end{document}
